\begin{table}\centering\tcp{\tsc{Seram-Tanimbar-Bomberai} *d = **d}{11-02}
\begin{adjustbox}{width=\textwidth}
	\begin{threeparttable}\st{2.5pt}
	\begin{tabular}{ll|llllll}\lsptoprule
*gloss	&		&	two	&	leaf	&	inside	&	blood	&	thorn\su{Δ}	&	stand\\
PMP	&	*d	&	*\tbr{d}uha	&	*\tbr{d}ahun	&	*\tbr{d}aləm	&	*\tbr{d}aRaq	&	*\tbr{d}uRi	&	*\tbr{d}iRi, *kə\tbr{d}əŋ\\	\midrule
Bobot	&	{\hr}l	&	{\hr}\tbr{l}ua	&	{\hr}ai \tbr{l}an	&	{\hr}vi \tbr{l}alan	&	{\hr}\tbr{l}awa	&	{\hr}-\tbr{l}ui-n	&	{\hr}a-ko\tbr{l}an\\
Masiwang	&	{\hr}l	&	{\hr}\tbr{l}ua	&	{\hr}\tbr{l}an	&	{\hr}ramuli \tbr{l}alan	&	{\hr}\tbr{l}asa	&	{\hr}\tbr{l}usi-n	&	{\hr}ko\tbr{l}an\\	\midrule
Bati	&	{\hr}l	&	‹\tbr{l}otu›	&	‹\tbr{l}ino›	&		&	‹\tbr{l}alai›	&	‹\tbr{l}ulul›	&	\\
Geser-G.	&	{\hr}r	&	{\hr}\tbr{r}oti	&	{\hr}\tbr{r}u	&	\cg{(lomin)}\su{†}	&	{\hr}\tbr{r}ara	&	{\hr}\tbr{r}uri	&	{\hr}ma-\tbr{r}iri\\
Watubela	&	{\hr}l	&	{\hr}\tbr{l}ua	&	{\hr}\tbr{l}ueʔ	&	\cg{(lomi)}\su{†}	&	{\hr}\tbr{l}alak	&	{\hr}\tbr{l}uli	&	{\hr}-\tbr{l}ili\\
Kesui	&	{\hr}r	&	‹\tbr{r}ua›	&	‹a\tbr{r}ehin›	&		&	‹\tbr{l}arah›	&	‹\tbr{l}úru›	&	\\
Kowiai	&	{\hr}r	&	{\hr}\tbr{r}uwet	&	{\hr}\tbr{r}owa	&	\cg{(mora)}\su{†}	&	{\hr}\tbr{r}ara	&	{\hr}\tbr{r}ur	&	{\hr}nama-\tbr{r}ir\\	\midrule
Teor	&	{\hr}r	&	‹\tbr{r}úa›	&	\cg{(‹chafen›)}	&		&	‹\tbr{l}arah›	&	\cg{(‹urut›)}	&	\\
Kur	&	{\hr}r	&	{\hr}\tbr{r}ua	&		&		&		&		&	\\	\midrule
Yamdena	&	{\hr}d	&	{\hr}\tbr{d}u	&	{\hr}\tbr{d}on-e	&	{\hr}\tbr{d}alam	&	{\hr}\tbr{d}ara-ny	&	{\hr}\tbr{d}uri	&	{\hr}nam\tbr{d}iry\\
Fordata	&	{\hr}r	&	{\hr}\tbr{r}ua	&	{\hr}\tbr{r}oa	&	{\hr}\tbr{r}ala-n	&	{\hr}\tbr{l}ara-n	&	{\hr}\tbr{l}uri-n	&	{\hr}n-\tbr{d}iri\\
Kei	&	{\hr}r	&	{\hr}\tbr{r}u	&	{\hr}\tbr{r}on	&	{\hr}\tbr{r}aa-n	&	{\hr}\tbr{l}ar	&	{\hr}\tbr{l}uri-n	&	{\hr}\tbr{d}ir\\
Onin	&	{\hr}n	&	{\hr}\tbr{n}ua	&	{\hr}wi \tbr{n}on-in	&	{\hr}\tbr{n}anam\su{‡}	&	{\hr}\tbr{r}ara	&	{\hr}\tbr{r}uri-r	&	{\hr}ma\tbr{r}iri\\
Sekar	&	{\hr}n	&	{\hr}\tbr{n}ua	&	\cg{(kafak-in)}	&	{\hr}\tbr{n}anam	&	{\hr}\tbr{r}ara	&	{\hr}\tbr{r}uri-r	&	{\hr}m\tbr{r}iri\\
Uruang.	&	{\hr}n	&	{\hr}\tbr{n}ua	&	{\hr}\tbr{n}on-in	&	{\hr}\tbr{n}anam	&	{\hr}\tbr{n}ara	&	{\hr}\tbr{n}uri-n	&	{\hr}a-m\tbr{r}iri\\
Arguni	&	{\hr}r	&	{\hr}\tbr{r}u	&	\cg{(náne)}	&	\cg{(ramo-n)}\su{†}	&	{\hr}\tbr{r}áre	&	\cg{(tór)}	&	\cg{(gáser)}\\
Bedoanas	&	{\hr}r	&	{\hr}\tbr{r}uyu	&	{\hr}\tbr{r}ína	&		&	\cg{(rita)}	&	\cg{(tóʔor)}	&	\cg{(usar)}\\
Erokwanas	&	{\hr}r	&	{\hr}\tbr{r}u	&	\cg{(manda-nánia)}	&	\cg{(-ramu)}\su{†}	&	\cg{(rita)}	&	\cg{(toʔar)}	&	\cg{(waser)}\\
Irarutu	&	{\hr}r	&	{\hr}\tbr{r}ifo	&	{\hr}\tbr{r}o	&	\cg{(gan)}	&	\cg{(wams)}	&	{\hr}\tbr{r}ur	&	{\hr}n-ma\tbr{r}ir\\
Kuri	&	{\hr}r	&	{\hr}\tbr{r}u	&	{\hr}\tbr{r}uɛˈrɔ	&	\cg{(gurˈnɛ)}	&	\cg{(wams)}	&	{\hr}\tbr{r}ur	&	{\hr}itə-b\tbr{r}ir\\
		\lspbottomrule
	\end{tabular}
		\begin{tablenotes}\tnsize
			\item[†]	{Despite superficial similarity,
								the Geser-Gorom, Kowiai, Watubela, Arguni, and Erokwanas words
								for ‘inside’ are not directly from *daləm.
								They instead reflect another form such as PCEMP **laman ‘deep’
								\citep{BlustTrussel2020} or a similar form,
								seen also in Yamdena \it{melaman} ‘deep’\it{,} Kei \it{laman}
								‘depth’, and Ile Mandiri Lamaholot \it{ləmeʔ} ‘deep’.
								(Kowiai \it{mora} has consonant metathesis from **laman.)}
			\item[‡]	{The Onin word list in \citet{Kaiping2019} has a typographical
								error with ‹naram› ‘inside’.
								The original word list on which this is based,
								\citet{Donohue2010}, contains medial \it{n}.
								Medial \it{n} is also confirmed by \citet[212]{SmitsVoorhoeve1992b}
								with Onin \it{ami nanam} ‘inside, within’.}
			\item[Δ]	{PMP *duRi is reconstructed with the meaning ‘thorn, splinter, fish bone’.
								The Yamdena, Sekar, and Onin reflexes mean ‘thorn,
								bone’, while other reflexes are only known to mean ‘bone’.}
		\end{tablenotes}
	\end{threeparttable}
\end{adjustbox}
\end{table}

\sscp{*d/*z > **d}{11-01-01}
The first change shared by STB is merger of *d/*z,
at least in word-initial position.
The data from Yamdena indicates that the outcome of this merger
was probably **d (discussed further below),
while other STB languages have reflexes which are likely reflexes of earlier **d.
Merger of *d/*z also occurs in \tsc{Ambon-Seram}
and in Banda, but in \tsc{Ambon-Seram} no language has *d/*z
> \it{d,} meaning that the outcome of the merger in Proto-Ambon-Seram
may have been different to Proto-STB.
See \srf{07-03-03} and \srf{13-04} for more discussion.
Examples of *d = **d are given in \trf{11-02}
and examples of *z > **d in \trf{11-03}. %Discussion follows.

\begin{table}\tcp{\tsc{Seram-Tanimbar-Bomberai} *z > **d}{11-03}
\begin{adjustbox}{width=\textwidth}
	\begin{threeparttable}\st{2.5pt}
	\begin{tabular}{ll|llllll}\lsptoprule
*gloss	&		&	path	&	far	&	rain	&	chin, jaw	&	ladder	&	bad, evil\\
PMP	&	*z	&	*\tbr{z}alan	&	*\tbr{z}auq	&	*qu\tbr{z}an	&	*qa\tbr{z}ay	&	*haRə\tbr{z}an	&	*\tbr{z}aqat\\	\midrule
Bobot	&	{\hr}l	&	{\hr}\tbr{l}olan	&	\cg{(otak)}	&	{\hr}u\tbr{l}an	&	{\hr}ya\tbr{l}a-n	&	{\hr}yo\tbr{l}an	&	\cg{(‹a′vet›)}\\
Masiwang	&	{\hr}l	&	{\hr}\tbr{l}olon	&	\cg{(wotat)}	&	\cg{(roti)}	&	{\hr}yala-n	&	{\hr}wo\tbr{l}an	&	\cg{(lalan sa)}\\	\midrule
Bati	&	{\hr}l	&	‹\tbr{l}āān›	&		&	‹u′an›	&		&		&	\\
Geser-G.	&	{\hr}r	&	{\hr}\tbr{l}alan	&	{\hr}\tbr{r}au	&	{\hr}u\tbr{r}an	&	{\hr}a\tbr{r}	&	{\hr}ro\tbr{r}an	&	\cg{(garata)}\\
Watubela	&	{\hr}l	&		&	{\hr}\tbr{l}au	&	\cg{(kadama)}	&	\cg{(kakelan)}	&	{\hr}la\tbr{l}on	&	\\
Kesui	&	{\hr}r	&	‹la\tbr{r}an› \it{(met.)}	&	\cg{(‹udáma›)}	&		&		&	‹\tbr{r}áhat›\\
Kowiai	&	{\hr}r	&	{\hr}\tbr{r}aran	&	{\hr}\tbr{s}or	&	\cg{(oma)}	&	{\hr}la\tbr{r}aʔai	&	{\hr}ro\tbr{r}an	&	\cg{(nuŋgutei)}\\	\midrule
Teor	&	{\hr}r	&	\cg{(‹lagain›)}	&		&	‹hu\tbr{r}ani›	&		&		&	‹\tbr{y}at›\\
Kur	&	{\hr}r	&	\cg{(liŋain)}	&		&	{\hr}u\tbr{r}an	&		&		&	‹\tbr{j}aaid› [\tbr{y}-]\\	\midrule
Yamdena	&	{\hr}d	&	{\hr}\tbr{d}alnen-e	&	{\hr}do{\tl}\tbr{d}o	&	{\hr}u\tbr{d}an	&	\cg{(keka-ny)}	&	\cg{(fnite)}	&	{\hr}\tbr{y}atak\\
Fordata	&	{\hr}r	&	\cg{(liŋaʔan)}	&	{\hr}ra{\tl}\tbr{r}oa	&	\cg{(daʔut)}	&	\cg{(demi-n)}	&	\cg{(fnita)}	&	\cg{(sian)}\\
Kei	&	{\hr}r	&	\cg{(laŋay)}	&	{\hr}ro{\tl}\tbr{r}o	&	\cg{(doʔot)}	&	\cg{(kakean)}	&		&	\cg{(sian)}\\
Onin	&	{\hr}n	&	{\hr}\tbr{n}aman	&	\cg{(bair)}	&	{\hr}u\tbr{n}an	&	\cg{(keken)}	&	\cg{(nene)}	&	\cg{(berik)}\\
Sekar	&	{\hr}n	&	{\hr}\tbr{n}anam	&	\cg{(bair)}	&	\cg{(yagin)}	&	\cg{(kaken)}	&	\cg{(neg)}	&	\cg{(kwes ta)}\\
Uruang.	&	{\hr}n	&	{\hr}\tbr{n}enmatan	&	\cg{(olat)}	&	{\hr}u\tbr{n}an	&	\cg{(keken)}	&	{\hr}\tbr{n}eran \cg{(met.)}	&	\cg{(tibyabin)}\\
Arguni	&	{\hr}r/d	&	{\hr}\tbr{r}árin	&	{\hr}\tbr{d}ówan	&	\cg{(úmun)}	&	\cg{(aure-m)}	&	\cg{(rumbur)}	&	\cg{(owi, ofi)}\\
Bedoanas	&	{\hr}r/d	&	{\hr}\tbr{r}árin	&	{\hr}\tbr{nd}ewen\su{†}	&	{\hr}á\tbr{r}i	&		&		&	{\hr}ó\tbr{y}et\\
Erokwanas	&	{\hr}r/d	&	{\hr}\tbr{r}órin	&	{\hr}\tbr{d}áwen	&	\cg{(úmin)}	&		&		&	{\hr}ú\tbr{y}et\\
Irarutu	&	{\hr}r	&	{\hr}\tbr{r}arn	&	{\hr}ne\tbr{r}o	&	\cg{(syem)}	&	\cg{(-kki-)}	&	{\hr}ron\su{‡}	&	\cg{(fit)}\\
Kuri	&	{\hr}r	&	{\hr}\tbr{r}arn	&	{\hr}gɛ\tbr{r}o	&	\cg{(nawnə̆)}	&	\cg{(ˈae)}	&	{\hr}rɔ\tbr{r}nɛ\su{‡}	&	\\
		\lspbottomrule
	\end{tabular}
		\begin{tablenotes}\tnsize
			\item[†]	{Bedoanas [nd] appears to be a variant of /d/,
								at least word initially.
								This is also attested in *dalij ‘buttress root’ > \it{dara {\tl} ndarao} ‘root’.
								(Arguni \it{dare(g)}, Erokwanas \it{dara} ‘root’).}
			\item[‡]	{The Irarutu
								and Kuri words for ‘ladder’ are from \citet[162]{SmitsVoorhoeve1992a},
								given as [rona] and [rɔrnɛ] respectively.
								Both languages have frequent vowel epenthesis after final consonants.
								Thus, these are analysable as /ron/ and /rorn/ respectively.
								Irarutu has *z > {\0} /{\gap}C while Kuri has *z > \it{r}.}
		\end{tablenotes}
	\end{threeparttable}
\end{adjustbox}
\end{table}

Discussion of some of the words in \trf{11-02} is needed.
Firstly, \tsc{Nuclear Tanimbar-Bomberai} (Fordata, Kei,
Onin, Sekar, Uruangnirin) have different reflexes of **d in the
environment /{\gap}V\it{r}, discussed in more detail in \srf{11-05-01}.
Kesui and Teor appear to have a similar process of regressive
dissimilation **r > \it{l} /{\gap}\it{Vr}.

Secondly, The changes affecting the initial consonant in *diRi
‘stand’ are also different,
being preserved as \it{d} in Fordata and Kei.
This is probably due to intermediate **nd/**md with the nasal
deriving from a stative prefix *ma-.
This prefix is still attested in several languages.

Thirdly, Yamdena has *d = \it{d} in word-initial position,
but word medially it has a split of *d > \it{r,
d}.\footnote{
		We thank a reviewer for bringing the reflexes of
		medial *d in Yamdena to our attention.
		For the sake of completeness the three word-final reflexes of
		*d identified are *tuhu\tbr{d} > \it{tu\tbr{r}-e} ‘knee’,
		*qabatə\tbr{d} > \it{bata\tbr{r}} ‘sago grub’,
		and *sumaŋə\tbr{d} > \it{smwaŋa\tbr{t}} ‘soul’.}
Five instances
of *d > \it{r} have been identified: *tuduq > \it{tu\tbr{r}u}
‘drip’,  *sida > \it{si\tbr{r}} ‘\tsc{3pl}’,
*maqudip > \it{mo\tbr{r}ip} ‘alive’,
*sudu > \it{su\tbr{r}w, su\tbr{r}u} ‘spoon’,
and  *udu > \it{u\tbr{r}u} ‘grass that grows about knee high’.
Three instances of *d = \it{d} have been identified: *maŋuda
> \it{maŋu\tbr{d}-e} ‘young (fruit)’,
*kudən > \it{ku\tbr{d}an} ‘pot’,
and *ma-udəhi > \it{mu\tbr{d}y} ‘behind, outside’.

Before we discuss specific words in \trf{11-03},
we must highlight that Yamdena may not have a merger of *d/*z
in word medial positions.
Word medially Yamdena has *d > \it{r, d} (discussed above).
Reflexes of medial *z in Yamdena are scarce with only three putative
reflexes identified: *quzan > \it{u\tbr{d}an} ‘rain’,
*quzuŋ > \it{u\tbr{nd}u-n} ‘top,
extremity, upper point’ \citep[115]{Drabbe1932b}, \it{u\tbr{nr}un}
‘tip, end, point’ \citep{MettlerMettler1996},\footnote{
		Yamdena ‹nd›
		in \citet{Drabbe1932b} is equivalent to /nr/ in
		\citet{MettlerMettler1996,MettlerMettler1997}.
		This phoneme is described by \citet{MettlerMettler1997} as following:
		“nr is most often pronounced [ndr] or [nᵈr],
		[nd] is also very common,
		mostly among younger people.
		Not very many people pronounce it as a pure [nr].
		However, the [r] sound is more dental,
		the tip of the tongue just behind the teeth.”} 
and perhaps *tuzuq > \it{-tu\tbr{t}u} ‘point at’.
Given these reflexes, and given *d > \it{r,
d}, it is possible to posit medial merger of *d/*z > **d in Yamdena,
with subsequent **d > \it{r} /V{\gap}V in most cases.
However, our confidence in this is much lower than it is for
word-initial *d/*z > \it{d}.

Discussion of some of the words in \trf{11-03} is also needed.
Firstly, regarding the reflexes of *zalan ‘path’
\citet{MettlerMettler1996} list multiple words for ‘road,
path’ in Yamdena: \it{dalamten-e,
dalnen-e, dalnyen-e,} and \it{dalŋen}.
These words are best seen as reflexes of PMP *zalan ‘road,
path’ with irregular changes to the final consonant in some cases,
rather than *zalan being replaced by **dalam (contra \citealt[276]{Blust1993}).
\citet[226]{Jackson2014} gives Irarutu \it{rarn} ‘path’ as being
from the Kaimana dialect of Irarutu.
He also gives variants \it{radni}
and \it{rarum} as from other (unspecified) dialects.
The Fruata dialect of Irarutu has \it{radni} [radəni] \citep{Voorhoeve1995b}.
Reflexes of *zalan ‘path’ in Sekar
and Onin have irregular final *n > \it{m}.
(Onin \it{naman} has consonant metathesis from earlier **nanam.)
Irregular final *n > \it{m} also occurs in the alternate Yamdena
form \it{dalamten-e}, as well as Irarutu \it{rarum}.
Liana-Seti (\tsc{Ambon-Seram}, \srf{10-03-04}) also has \it{lalam} alongside \it{lala}.
Geser-Gorom \it{lalan} ‘path’ has regressive assimilation of
**r > \it{l} /{\gap}Vl.
Kesui \it{laran} ‘path’ probably has consonant metathesis from earlier **ralan.

Regarding putative reflexes of *zauq,
Kowiai \it{sor} may be a reflex with *z > \it{s} rather than
*z > \it{r}, though the origin of final \it{r} is unknown.
\citet[112]{Jackson2014} gives Irarutu \it{ne-iro} as a possible
analysis of \it{nero} ‘far’.
The basis for this analysis is unexplained
and any meaning for putative \it{ne-} is unknown.
Nonetheless, the Kuri word with initial \it{ge} may attest a
different (historical) prefix and indicate that only the
final syllable of these words is inherited from *zauq.

Bati has *d/*z > **l > \it{l,} {\0} /V{\gap}V,
thus accounting for *z > {\0} in *quzan > ‹u′an›.
The symbol ‹′› in \citet{Wallace1869} is of unclear value.
It may represent a glottal stop,
but also frequently occurs in contexts where no glottal stop
occurs, such as Malay ‹a′njing› [ˈaɲʤiŋ] ‘dog’.
Other Bati words with *l > {\0} in the available data are: *zalan
> ‹\it{lāān}› ‘road’, *bulan > ‹wúan› ‘moon’,
as well as **kazupay > ‹kaúfei› ‘mouse,
rat’ with *z > **l > {\0} (cf.
Geser \it{kalufa} [\citealp[144]{Collins1986}]).

\tsc{Teor-Kur} and \tsc{Tanimbar-Bomberai} have *z > \it{y} in reflexes of *zaqat.
This may be a lexically specific innovation shared between these
two branches of STB.
Yamdena \it{yatak} could be linked to *zaqat but the pathway
between the two requires positing several irregular sound changes.
It could involve consonant metathesis to **zataq with irregular
*q > \it{k}, or the entire final syllable may be a historical suffix,
though both are ad-hoc hypotheses.
The origin of the initial vowels in Bedoanas \it{óyet}
and Erokwanas \it{úyet} are also not explained if these are inherited from *zaqat.

Apart from the reflexes of *z given in \trf{11-03},
reflexes of PCEMP **kazupay ‘mouse,
rat’ also occur in STB languages: Geser \it{kalufa} \citep[144]{Collins1986},
\it{karufa} \citep{Stokhof1982b}, Gorom \it{aruha} \citep{Visser2019a},
Bati ‹kaúfei›, Kesui \it{arofa,} Kowiai \it{asafai},
and Irarutu \it{sfe} [səfe].
These reflexes have various irregularities indicating that they spread by contact.
The irregularities are: *z > \it{l {\tl} r} in Geser (expect
*z > \it{r}), *ay > ‹ei› in Bati (expect *ay > \it{e}),
*z > s in Irarutu,
as well as *z > \it{s}
and *ay > \it{ai} in Kowiai (expect *z > \it{r}
and *ay > \it{a}).
The Sera dialect of Fordata has \it{kandruə} ‘rat’ which may
be from **kazupay or **kandoRa ‘cuscus’.

